\chapter{A radial obstacle problem and its contact transition}\label{sec:examples-radial-obstacle}
We solve the normalized obstacle problem in a ball, with obstacle zero and constant boundary height. The normalization is \cite{Fig18}, equation (5.2), p.~9, and Definition \ref{def:normalized-obstacle}; the explicit radial computations below illustrate that equation and the regular and singular models of Sections 7.1--7.2 of the source.

\begin{proposition}[A radial profile with a prescribed contact ball]\label{prop:radial-obstacle-profile}
Let $n\geq2$ and $a>0$. Put $U_a(r)=0$ for $0\leq r\leq a$ and, for $r>a$, set
\[
 U_a(r)=\int_a^r\frac{s^n-a^n}{ns^{n-1}}\,ds.
\]
Then $u_a(x)=U_a(|x|)$ is a globally $C^{1,1}$ normalized obstacle solution on $\mathbb R^n$, with contact set $\overline B_a$ and positivity outside that ball. Explicitly,
\[
 U_a(r)=\begin{cases}
 \dfrac{r^2-a^2}{4}-\dfrac{a^2}{2}\log(r/a),&n=2,\\[4pt]
 \dfrac{r^2-a^2}{2n}+\dfrac{a^n}{n(n-2)}(r^{2-n}-a^{2-n}),&n>2
 \end{cases}\qquad(r>a).
\]
\end{proposition}
\begin{proof}
The radial equation is $(r^{n-1}U')'=r^{n-1}$. Integrating with $U'(a)=0$ gives $U'(r)=(r^n-a^n)/(nr^{n-1})$, and a second integration with $U(a)=0$ gives the stated formulas. For $r>a$ the derivative is positive, so the contact set is exactly the closed ball. The tangential and radial eigenvalues of its Hessian there are
\[
 \frac{U_a'}r=\frac1n\left(1-(a/r)^n\right),\qquad
 U_a''=\frac1n+\frac{n-1}{n}(a/r)^n.
\]
They lie in $[0,1]$ and have sum, with tangential multiplicity $n-1$, equal to one. On the contact ball the gradient and Hessian vanish. Since $U_a(a)=U_a'(a)=0$, the gradient is continuous across the sphere and no surface measure appears in its distributional derivative. Its bounded weak Hessian makes the gradient globally Lipschitz, with bound one. The function is therefore $C^{1,1}$ and has Laplacian equal to one outside contact and zero inside, as required. It is not $C^2$ across the sphere: the radial second derivative jumps from zero to one.
\end{proof}

\begin{theorem}[All constant-height Dirichlet cases in a ball]\label{thm:radial-obstacle-dirichlet}
Fix $n\geq2$, $R>0$ and $h\geq0$. The unique minimizer of
\[
 J(v)=\int_{B_R}\left(\tfrac12|Dv|^2+v\right),\qquad
 v\geq0,\quad v-h\in H^1_0(B_R),
\]
is as follows, with $h_*=R^2/(2n)$:
\[
 u_h(x)=\begin{cases}
 0,&h=0,\\
 U_a(|x|),\quad U_a(R)=h,\ 0<a<R,&0<h<h_*,\\
 |x|^2/(2n),&h=h_*,\\
 h-h_*+|x|^2/(2n),&h>h_*.
 \end{cases}
\]
In the middle range the contact radius $a$ is uniquely determined and decreases strictly as $h$ increases.
These explicit solutions illustrate the normalized obstacle equation in \cite{Fig18}; the proof below verifies minimality against all admissible competitors.
\end{theorem}
\begin{proof}
Existence and uniqueness follow from Theorem \ref{thm:obstacle-existence} with forcing $f=-1$. For $0<a<R$ set $g(a)=U_a(R)$. Differentiating its integral expression, whose integrand vanishes at the moving lower endpoint, gives
\[
 g'(a)=-a^{n-1}\int_a^R s^{1-n}\,ds<0.
\]
The explicit formulas show $g(a)\to h_*$ as $a\downarrow0$ and $g(a)\to0$ as $a\uparrow R$. Continuity and strict monotonicity supply exactly one radius for each intermediate $h$.

Every displayed function is nonnegative, has boundary value $h$, and is $C^{1,1}$ with $\Delta u_h=\chi_{\{u_h>0\}}$ almost everywhere. For a competitor $v$, integration by parts yields
\[
 \int_{B_R}Du_h\cdot D(v-u_h)+\int_{B_R}(v-u_h)
 =\int_{\{u_h=0\}\cap B_R}v\geq0.
\]
Thus each candidate satisfies the variational inequality and is the unique minimizer. This verifies the global boundary-value problem directly; rotational symmetry alone is not used to assume away possible nonradial competitors.
\end{proof}

\begin{proposition}[Contact is not the same as positive reaction mass]\label{prop:radial-obstacle-reaction}
For $0<h<h_*$ the reaction measure of $u_h$ is $\lambda_h=\chi_{B_a}\mathcal L^n$ and its mass is $\omega_n a^n$. At $h=h_*$ the contact set is the singleton $\{0\}$ but the reaction is zero. At $h>h_*$ contact is empty and the reaction remains zero; at $h=0$ the whole ball is in contact and $\lambda_0=\mathcal L^n\lfloor B_R$.
\end{proposition}
\begin{proof}
With forcing $f=-1$, Proposition \ref{prop:obstacle-reaction} gives $\lambda_h=1-\Delta u_h$. Substitute the Laplacians computed in Theorem \ref{thm:radial-obstacle-dirichlet}. The sphere $\partial B_a$ has zero Lebesgue measure and carries no singular reaction, since the gradient is continuous across it. At the threshold the Laplacian of $|x|^2/(2n)$ is one even at zero, while its value vanishes there. Thus a nonempty contact set need not carry nonzero reaction measure.
\end{proof}

\begin{proposition}[The regular sphere and the singular point have different blow-ups]\label{prop:radial-obstacle-blowups}
For $0<h<h_*$ and $p=ae\in\partial B_a$, $|e|=1$, the full family of quadratic rescalings satisfies
\[
 \frac{u_h(p+ry)}{r^2}\longrightarrow\frac12((y\cdot e)^+)^2
 \quad\text{in }C^1_{\mathrm{loc}}(\mathbb R^n)\quad(r\downarrow0).
\]
At $h=h_*$ the blow-up at the sole contact point is exactly $|y|^2/(2n)$ at every scale. Hence the contact sphere consists of regular points, whereas the threshold point is singular with a positive-definite blow-up matrix.
\end{proposition}
\begin{proof}
Near $a$, Taylor's theorem on the exterior side gives $U_a(a+t)=\frac12(t^+)^2+O(|t|^3)$ and $U_a'(a+t)=t^++O(t^2)$, with zero values for $t\leq0$. On bounded $y$ sets,
\[
 |p+ry|-a=r(y\cdot e)+O(r^2),\qquad
 \frac{p+ry}{|p+ry|}=e+O(r).
\]
Substitution into the function and radial gradient formulas gives uniform convergence of both the rescaled function and gradient. The estimates remain valid across the changing positivity interface because $t\mapsto t^+$ is Lipschitz. At the threshold, exact quadratic homogeneity gives the second assertion. Apply Definitions \ref{def:obstacle-regular-point} and \ref{def:obstacle-singular-point}; the singular matrix is $I/n$ and has zero-dimensional kernel.
\end{proof}

\begin{proposition}[How the contact radius approaches its endpoints]\label{prop:radial-obstacle-contact-asymptotics}
As $h\downarrow0$, the contact radius satisfies $R-a\sim\sqrt{2h}$. As $h\uparrow h_*$, it tends to zero. In dimension two the exact deficit is
\[
 h_*-h=\frac{a^2}{2}\log(R/a)+\frac{a^2}{4}.
\]
In dimensions $n>2$ it is
\[
 h_*-h=\frac{a^2}{2(n-2)}-\frac{a^nR^{2-n}}{n(n-2)},
 \qquad a\sim\sqrt{2(n-2)(h_*-h)}.
\]
Here $f\sim g$ means $f/g\to1$ in the indicated limit.
\end{proposition}
\begin{proof}
Subtract the explicit value $U_a(R)$ from $h_*$ to obtain the deficit identities. For $n>2$ the second term is lower order than $a^2$ as $a\downarrow0$, giving the stated asymptotic. In dimension two the logarithm grows without bound and cannot be dropped. At the other endpoint, differentiate the formula for $g'(a)$ in Theorem \ref{thm:radial-obstacle-dirichlet}; it gives $g(R)=g'(R)=0$ and $g''(R)=1$ by continuous extension from the left. Taylor's theorem yields $h=g(a)=(R-a)^2/2+o((R-a)^2)$, proving the first assertion. These calculations exhibit a change in contact geometry even though the minimizer remains controlled in $C^{1,1}$.
\end{proof}

\begin{proposition}[The derivative of the minimized energy]\label{prop:radial-obstacle-energy-response}
Let $E(h)=J(u_h)$ for $h\geq0$. It is continuously differentiable from the right at zero and differentiable elsewhere, with
\[
 E'(h)=\begin{cases}
 \omega_n(R^n-a(h)^n),&0<h<h_*,\\
 \omega_nR^n,&h\geq h_*,
 \end{cases}
 \qquad E'(0+)=0.
\]
It is strictly convex on $(0,h_*)$ and affine on $[h_*,\infty)$.
\end{proposition}
\begin{proof}
In the intermediate range, $g'(a)<0$ makes $a(h)$ smooth. Differentiate the energy as a radial integral over $a<|x|<R$. There is no moving-interface contribution because both $U_a$ and its gradient vanish at $r=a$. Integration by parts in the resulting first variation leaves only the outer boundary term: the interior term has coefficient $1-\Delta u_h=0$ on positivity, and the parameter derivative is zero in contact. The boundary derivative with respect to $h$ is one. Consequently
\[
 E'(h)=\mathcal H^{n-1}(\partial B_R)U_a'(R)
 =n\omega_nR^{n-1}\frac{R^n-a^n}{nR^{n-1}}
 =\omega_n(R^n-a^n).
\]
This derivative is strictly increasing because $a'(h)<0$. Above threshold the solution is the threshold solution plus the constant $h-h_*$, so its gradient energy is unchanged and its linear energy increases by $(h-h_*)|B_R|$.

The explicit profiles and their gradients converge uniformly on the closed ball at both limiting heights, so $E$ is continuous there. The interior derivative has limits zero at $h=0$ and $\omega_nR^n$ at $h=h_*$. The mean value theorem, applied on one-sided intervals and then taking their endpoints to the threshold, gives the claimed matching derivatives. This proves the endpoint assertions without differentiating an undefined inverse radius at zero.
\end{proof}
